\begin{abstract}
Small scale liquid rocket engines are seeing renewed interest for satellite propulsion, lunar landers and amateur liquid rockets, but pump fed feed systems at this scale are poorly documented and the open data that exists is almost entirely computational. A small engine needs high pressure at low flow, which drives the pump to a low specific speed where the radial impeller bleeds efficiency to disc friction, and drives the turbine that powers it to a low blade to jet speed ratio with partial admission. This thesis investigates a 20,000 rpm turbopump built around a low specific speed Barske partial emission pump ($n_q$ of 6.4, flow rate of 0.3 L/s, 20 bar rise on water) coupled to a supersonic partial admission impulse turbine (2.6 kW, 7.65 bar inlet, on compressed air), to supply the experimental characterisation this regime lacks.

The pump was sized from Barske's and Lock's models and the turbine from a method of characteristics blade with the Goldman starting and Sasman and Cresci separation criteria and the Weiss loss model. Both were built, run on a custom data acquisition system, and tested under electric motor drive and as a coupled unit, with a transparent casing that let the cavitation be filmed and frame matched to the data.

The pump delivers its design head, with a head coefficient of 1.10 ± 0.05 collapsing across five speeds to confirm dynamic similarity. Peak measured efficiency was (19 ± 2)\%, dominated by a disc friction churning loss measured at 2.5 to 2.9 times Barske's estimate, extrapolating to about 26\% at the design speed. The cavitation breakdown was shown, both in the data and on video, to be set by the diffuser throat rather than the impeller eye, confirming Lock's diffuser limited picture. Coupled operation was demonstrated at around 5400 rpm, and the measured torque deficit was explained by a predicted rotor unstart, confirming the Goldman starting criterion off design. The models used are sufficient for design stage prediction of this class of machine.
\end{abstract}
